% !TeX program = xelatex
% UTF-8 transcription of all 49 source photographs, in filename/page order.
% Original JPEG bytes are preserved in images/pic01.jpg through images/pic49.jpg.
% Source wording, abbreviations and doses are transcribed, not medically corrected.
% Repeated printed stationery labels: موضوع ؛ تاریخ / / .
% Blank date fields and ruled lines have no textual content.
% Compile this file directly with XeLaTeX, or input it into a XePersian document.
\ifdefined\settextfont
  \def\ViralInfectionsStandalone{0}
\else
  \def\ViralInfectionsStandalone{1}
  \documentclass[12pt,a4paper]{article}
  \usepackage[margin=2cm]{geometry}
  \usepackage{amsmath,amssymb,xcolor,enumitem,graphicx,array,longtable}
  \usepackage{xepersian}
  \IfFontExistsTF{Amiri}{\settextfont{Amiri}}{\settextfont{FreeSerif}}
  \setlatintextfont{FreeSerif}
  \setlength{\emergencystretch}{2em}
\fi
\providecommand{\en}[1]{\lr{#1}}
\newcommand{\vih}[1]{\colorbox{yellow!35}{\strut #1}}
\newcommand{\viunclear}[1]{\textcolor{red!65!black}{[خوانش نامطمئن: #1]}}
\newcommand{\vipage}[3]{\clearpage\section*{تصویر #1 ـ صفحهٔ دست‌نویس #2}%
  {\footnotesize\en{#3}\par}}
\ifnum\ViralInfectionsStandalone=1 \begin{document}\fi

\vipage{۱}{۱}{pic01.jpg — 20261016\_043401147.jpg}
\section*{تظاهرات پوستی عفونت‌های ویروسی ـ ۲۵}
\begin{latin}Rook's Dermatology (10th)\end{latin}
\subsection*{هرپس‌ویروس‌ها}
\begin{itemize}
\item \en{DNA} ویروس هستند.
\item پس از بهبودی این ویروس‌ها از بدن پاک نمی‌شوند و در نورون‌ها می‌مانند و ممکن است باعث عفونت \en{latent} ایجاد کنند.
\end{itemize}
\subsection*{هرپس سیمپلکس \en{(HSV) (HHV-1)}}
\begin{itemize}
\item ویروس هرپس انسان.
\item از شایع‌ترین عفونت‌های ویروسی در انسان است.
\item ۲ تیپ آنتی‌ژنی شایع دارد:
\begin{itemize}
\item \vih{\en{Type 1}}: عفونت \vih{\en{facial (oral)}}.
\item \vih{\en{Type 2}}: عفونت \vih{ژنیتال}.
\item البته که \en{overlap} وجود دارد.
\end{itemize}
\item پس از عفونت اولیه ویروس \en{latent} می‌شود.
\end{itemize}
\subsubsection*{عفونت اولیه}
\begin{itemize}
\item معمولاً \vih{\en{subclinical}} است اما اگر تظاهرات پوستی در آن رخ دهند شدیدتر از \en{recurrence} است.
\item عفونت اولیهٔ ژنیتال با احتمال بالاتری علامت‌دار می‌شود.
\end{itemize}
\subsubsection*{عود عفونت}
\begin{itemize}
\item می‌تواند علامت‌دار یا تحت‌بالینی باشد.
\item احتمال عود: تیپ یک ۸۰\% و تیپ دو ۹۰\% است.
\item تحریک‌کننده‌های عود: تروما، عفونت‌های تب‌دار، \en{URI}، اشعهٔ \en{UV}، نورالژی تری‌ژمینال، جراحی‌های اینتراکرانیال، جراحی‌های دنتال، پروسیجرهای زیبایی صورت، لیزر.
\end{itemize}

\vipage{۲}{۲}{pic02.jpg — 20261016\_043407498.jpg}
\begin{itemize}
\item احتمال عود در زنان در پره‌منسترال بیشتر است.
\item استرس روانی تحریک‌کنندهٔ عود می‌باشد.
\item ممکن است علت عود یافت نشود.
\item در محل عود، وزیکول‌ها کوچک هستند؛ به صورت گروهی کنار هم هستند.
\item علائم سیستمیک نادرتر از عفونت اولیه است.
\item \vih{طبق یک قانون}، در افراد با ایمنی سالم، وزیکول‌ها هرگز مخاط بوکال را درگیر نمی‌کنند.
\end{itemize}
\subsubsection*{شدینگ \en{(shedding)} ساب‌کلینیکال ویروس}
\begin{itemize}
\item در تیپ ۲ شایع‌تر است.
\item افرادی که اپیزودهای عود داشته‌اند در ۲۰\% روزها و افراد بدون عود در ۱۰\% روزها شدینگ دارند.
\item در سال اول بیشتر از سال‌های بعد است.
\end{itemize}
\subsubsection*{بیولوژی عفونت ویروس}
\begin{itemize}
\item \vih{انتقال از طریق تماس مستقیم با} \vih{ترشحات و یا بزاق}.
\item انتقال به گانگلیون‌های نورون‌های حسی.
\item \vih{در فاز نهفته قابل شناسایی نیست چون} مدام تولید نمی‌شود.
\item شدینگ در مخاط و نزدیکی ژنیتال.
\item در هر نقطه‌ای از پوست می‌تواند درگیری ایجاد کند.
\end{itemize}

\vipage{۳}{۳}{pic03.jpg — 20261016\_043412034.jpg}
\begin{itemize}
\item خودتلقیحی به انگشت ـ به دلیل گاز گرفتن یا مکیدن انگشت.
\item هر دو ایمنی سلولی و هومورال نقش دارند ولی نقش ایمنی سلولی بسیار مهم‌تر است.
\item سیستم ایمنی نمی‌تواند ایمنی کامل در مقابل \en{HSV} ایجاد کند.
\item درگیری نقطه‌ای که قبلاً درگیر نشده در بیمار با سابقهٔ عود را \en{non-primary first-episode infection} می‌گویند.
\item بیماری که قبلاً فرم دهانی را گرفته اگر به فرم ژنیتال مبتلا شود، بیماری خفیف‌تر خواهد بود.
\item \viunclear{واژهٔ اول این سطر}، نقص ایمنی، \viunclear{واژهٔ پیش از «داروها»} داروهای موضعی احتمال عود را بالا می‌برند.
\item احتمال عود در بیماران درماتیت آتوپیک بیشتر است.
\end{itemize}
\subsection*{\en{HSV} که دهان و لب را درگیر کند، ۳ نام دارد}
\begin{itemize}
\item پرایمری هرپتیک ژنژیواستوماتیت.
\item هرپس لابیالیس.
\item \en{Cold sore}.
\end{itemize}
\begin{itemize}
\item شایع‌ترین تظاهر بالینی \en{HSV-1} است.
\item \vih{بیشتر در شیرخواران و کودکان کوچک رخ می‌دهد}.
\item معمولاً ساب‌کلینیکال و خفیف هستند.
\item به طور نادر ممکن است \vih{استوماتیت وزیکولار} ایجاد کند که بیمار دچار \viunclear{ادامهٔ عبارت در حاشیهٔ چپ تصویر بریده است}.
\end{itemize}

\vipage{۴}{۴}{pic04.jpg — 20261016\_043418918.jpg}
\begin{itemize}
\item بیشتر در کودکان \vih{\en{1--5}} سال رخ می‌دهد.
\item دورهٔ کمون حدود \vih{۵ روز}.
\item تب شدید، ضعف، آبریزش، بلع دردناک، بوی بد دهان، تورم لثه‌ها؛ لثه به سادگی خونریزی می‌کند.
\item وزیکول‌ها به صورت ضایعات پلاک سفیدرنگ در زبان، حلق، کام، مخاط بوکال.
\item \vih{سردرد می‌تواند بروز کند}.
\item تورم و تندرنس لنف‌نودهای منطقه‌ای.
\item \en{DDX}: استرپتوکوک، دیفتری، آفت، کاندیدا، بهجت.
\item \en{Complication}: فارنژیت در ۱۰\% موارد هرپس اوروژنیتال.
\begin{itemize}
\item فرم منتشر: نوزادان، نقص ایمنی.
\item در نوزادها با تظاهرات پوستی منتشر همراه است.
\item انسفالیت: در هر سنی کشنده؛ هپاتیت آن نادر است ولی می‌تواند کشنده باشد.
\end{itemize}
\item سیر بیماری: ظرف \vih{\en{3--5}} روز تب رفع می‌شود؛ در ظرف ۲ هفته بیماری کاملاً خوب می‌شود.
\end{itemize}
\subsubsection*{بررسی‌ها}
\begin{itemize}
\item سواب حاوی مایع وزیکول: \en{PCR}، کشت.
\item آنتی‌بادی: \vih{صرفاً سروکانورژن ملاک است}؛ به درد عودها نمی‌خورد.
\item بررسی \en{CSF}.
\end{itemize}

\vipage{۵}{۵}{pic05.jpg — 20261016\_043422205.jpg}
\subsubsection*{درمان}
\begin{itemize}
\item فرم خفیف بدون عارضه نیاز به درمان ندارد.
\item ضد درد به صورت دهان‌شویه.
\item آنتی‌سپتیک‌های موضعی برای جلوگیری از عفونت ثانویهٔ باکتریایی.
\item آنتی‌وایرال طول دورهٔ بیماری را کم می‌کند.
\end{itemize}
\paragraph{خط اول برای \en{less severe} که بلع مختل نیست}
\begin{itemize}
\item آسیکلوویر خوراکی \en{200 mg}، ۵ بار در روز، برای \vih{۵ روز}.
\item آسیکلوویر \en{800 mg BD}، برای ۵ روز.
\item سوسپانسیون آسیکلوویر \en{15 mg/kg}، ۵ بار، برای \vih{۷ روز}.
\item والاسیکلوویر خوراکی \en{500 mg BD}، برای ۵ روز.
\item فامسیکلوویر \en{250 mg BD}، برای ۷ روز.
\end{itemize}
\paragraph{خط دوم برای فرم شدید، نقص ایمنی}
\begin{itemize}
\item شروع \en{ASAP} درمان.
\item آسیکلوویر وریدی \en{5 mg/kg q8h}.
\item اگر انسفالیت: \en{10 mg/kg}.
\item افزایش \en{Cr} ممکن است؛ تجویز طی ۱ \viunclear{باقی عبارت در حاشیهٔ چپ بریده است}.
\end{itemize}

\vipage{۶}{۶}{pic06.jpg — 20261016\_043428165.jpg}
\subsection*{ریکارنت ژنژیواستوماتیت}
\begin{itemize}
\item ممکن است علامت‌دار یا تحت حاد باشد (صرفاً شدینگ).
\item احساس \vih{خارش یا سوزش}؛ پس از \vih{\en{1--2}} ساعت، وزیکول‌های کوچک دسته‌ای در یک زمینهٔ ملتهب.
\item درد می‌تواند از پیش بروز یا تا چند روز ادامه داشته باشد.
\item در کودکان وزیکول‌ها ممکن است بزرگ باشند.
\item فرم وزیکول‌ها: اغلب \en{irregular}، \en{in line}، \en{zosteriform}؛ با درد عمقی و لنفادنوپاتی همراه است.
\item در مرحلهٔ درگیری پوست و مخاط اغلب تب، درد، لنفادنیت وجود دارد.
\item در بعضی موارد ممکن است درد \vih{نورالژیک} \en{1--3} روز قبل شروع شود:
\begin{itemize}
\item سندرم \vih{\en{Mauriac}}.
\item همانند نورالژی تری‌ژمینال.
\end{itemize}
\item در موارد نقص ایمنی ممکن است اولسرهای ماندگار ایجاد شود.
\item درگیری چشم: شاخیهٔ چشم \viunclear{واژهٔ پیش از «چشم»}.
\item حملات متعدد در یک اندام، لنف‌ادم ایجاد می‌کند.
\item ظرف \vih{\en{7--10}} روز بدون اسکار خوب می‌شود.
\end{itemize}

\vipage{۷}{۷}{pic07.jpg — 20261016\_043431779.jpg}
\subsection*{اریتم مولتی‌فرم مرتبط با هرپس}
\begin{itemize}
\item ۹۵\% بیماران \en{EM} راجعه، سابقهٔ هرپس لابیالیس دارند.
\item هرپس \vih{چند روز تا ۲ هفته} قبل از \en{EM} رخ می‌دهد.
\item \en{HSV Ag gB} در ضایعات \en{EM}.
\item اگر درمان هرپس \vih{ظرف ۵ روز اول} با آسیکلوویر شروع شود از بروز \en{EM} جلوگیری می‌کند.
\item در حملات مکرر، دورهٔ ۶ ماهه آسیکلوویر پروفیلاکسی.
\end{itemize}
\subsection*{فلج بل \en{(Bell palsy)} مرتبط با هرپس \en{(HSV, VZV)}}
\begin{itemize}
\item درمان ضد ویروسی.
\item \en{PCR} بزاق باید انجام شود.
\end{itemize}
\subsection*{مننژیت لنفوسیتیک راجعه}
\begin{itemize}
\item در زمینهٔ \en{HSV} است.
\item یک مننژیت آسپتیک خوش‌خیم.
\item \en{3--14} روز طول می‌کشد.
\item انتروال‌های ماهانه تا سالانه.
\item آسیکلوویر پروفیلاکسی.
\end{itemize}

\vipage{۸}{۸}{pic08.jpg — 20261016\_043437443.jpg}
% Crossed-out heading in source: اریتم مولتی‌فرم.
\subsubsection*{سیر و پروگنوز}
\begin{itemize}
\item ضایعات پوستی از وزیکول تبدیل به پوستول شده، دلمه‌بندی می‌شوند و در نهایت ظرف \vih{\en{7--10} روز} بدون اسکار از بین می‌روند.
\end{itemize}
\subsection*{منجمنت هرپس لابیالیس راجعه}
\begin{itemize}
\item اگر حمله خفیف و کم‌فرکانس باشد، نیاز به درمان ندارد.
\item درمان: آسیکلوویر، والاسیکلوویر، فامسیکلوویر.
\item در افراد ایمونوکامپتنت مقاومت به آسیکلوویر نداریم.
\item در افراد با نقص ایمنی که متعدد آسیکلوویر گرفته‌اند مقاومت هرپس به آسیکلوویر مشاهده شده است؛ فوسکارنت یا سیدوفوویر بدهیم.
\item استفاده از \vih{ضد آفتاب} احتمال عود هرپس را کم می‌کند.
\item \vih{قبل از لیزر} زیبایی صورت، پروفیلاکسی هرپس بدهیم.
\end{itemize}
\paragraph{خط اول}
\begin{itemize}
\item کرم مرطوب‌کننده.
\item درسینگ هیدروکلوئید.
\item احساس سوزش را از بین می‌برند ولی در بهبودی اثری ندارند. % Marginal brace refers to the preceding two items.
\item \vih{پنسیکلوویر موضعی}: هر \en{2 h} برای ۴ روز.
\item آسیکلوویر خوراکی: \vih{\en{ASAP}}، \en{200 mg}، ۵ بار، ۵ روز.
\item والاسیکلوویر، فامسیکلوویر: \en{2 g BD} یک روز؛ یا \en{1500 mg} یک بار.
\end{itemize}

\vipage{۹}{۹}{pic09.jpg — 20261016\_043441279.jpg}
\begin{itemize}
\item تری‌فلورو تیمیدین: اثرش اثبات‌شده نیست.
\item ایدوکسوریدین موضعی ایمن نیست.
\end{itemize}
\paragraph{خط دوم (اگر \vih{عود مکرر} رخ داد)}
\begin{itemize}
\item آسیکلوویر پروفیلاکسی طولانی‌مدت: \en{200--400 mg BD for 4--6 months}.
\item والاسیکلوویر برای \en{6--12} ماه:
\begin{itemize}
\item بالای ۱۰ بار در سال: \en{1 g} روزانه.
\item زیر ۱۰ بار در سال: \en{500 mg} روزانه.
\end{itemize}
\end{itemize}
\paragraph{خط سوم}
\begin{itemize}
\item بیمار نقص ایمنی.
\item آسیکلوویر وریدی.
\item پنسیکلوویر.
\item \en{PEP} در بیمارانی که می‌خواهیم زیر تیغ ببریم:
\begin{itemize}
\item آسیکلوویر \en{PO} یا \en{IV}.
\item از چند روز قبل شروع کنیم.
\end{itemize}
\item موارد شدید مقاوم به آسیکلوویر: فوسکارنت، سیدوفوویر.
\end{itemize}

\vipage{۱۰}{۱۰}{pic10.jpg — 20261016\_043451522.jpg}
\subsection*{پرایمری هرپس ژنیتالیس}
\begin{itemize}
\item هر دو ویروس \en{HSV-1} و \en{HSV-2} می‌توانند ولی اغلب \en{HSV-2}.
\item انتقال عموماً جنسی است.
\item اکثراً پس از بلوغ؛ در کودکان سوءاستفادهٔ جنسی بررسی شود.
\item حتی در افراد بدون علامت، \vih{شدینگ در ۱۰\% روزها} رخ می‌دهد.
\end{itemize}
\subsubsection*{تظاهرات بالینی}
\begin{itemize}
\item \en{$\tfrac{1}{3}$} موارد علامت‌دار می‌شوند.
\item \vih{دورهٔ کمون}: ۲ روز تا ۲ هفته.
\item شایع‌ترین تظاهر، اروژن‌های دردناک است؛ خارش ضایعات، وزیکول و ادم.
\item شایع‌ترین محل ضایعات در مردان: روی گلنس و بدنهٔ آن، شفت است.
\item در زنان: اکسترنال ژنیتالیا \en{$<$} مخاط وولو \en{$<$} واژن \en{$<$} سرویکس.
\item در سیاه‌پوستان: \en{Post inflammatory hypopigmentation}.
\item پری‌آنال نیز دیده می‌شود.
\item عفونت سرویکس می‌تواند منجر به \en{Ulcerative Cervicitis} شود.
\item \vih{دیزوری} شایع است در هر دو جنس.
\item ممکن است ترشح \vih{مجرا} وجود داشته باشد.
\item لنفادنیت شایع است.
\end{itemize}

\vipage{۱۱}{۱۱}{pic11.jpg — 20261016\_043454443.jpg}
\begin{itemize}
\item در بیماران \en{HIV} ممکن است زخم‌ها مزمن شده و شبیه به ضایعات نئوپلاسمی باشند.
\end{itemize}
\subsubsection*{عوارض هرپس ژنیتال اولیه}
\begin{itemize}
\item رادیکولونوروپاتی: پارستزی ساکرال، رتنشن ادراری، یبوست، ایمپوتنسی (مردان)؛ ظرف چند روز تا چند هفته رفع می‌شود.
\item سردرد و مننژیسم: در ۳۶\% زنان و ۱۱\% مردان.
\item \en{PCR} مایع \en{CSF} مثبت است؛ لازم به \viunclear{ادامهٔ عبارت در حاشیهٔ چپ بریده است}.
\item انسفالیت نادر است؛ \en{full recovery}.
\end{itemize}
\subsubsection*{سیر بیماری و پروگنوز}
\begin{itemize}
\item ضایعات اگر درمان نشوند \vih{\en{2--3} هفته} باقی می‌مانند.
\item عود شایع است خصوصاً در \en{HSV-2}.
\item بررسی‌ها مانند پرایمری اوروژنیتال است.
\end{itemize}
\subsubsection*{منجمنت}
\begin{itemize}
\item استفاده از \en{barrier contraception}.
\item \en{Local microbicides}.
\end{itemize}

\vipage{۱۲}{۱۲}{pic12.jpg — 20261016\_043500796.jpg}
\paragraph{خط اول}
\begin{itemize}
\item آسیکلوویر خوراکی، والاسیکلوویر، فامسیکلوویر.
\item \vih{کنترل درد}.
\end{itemize}
\paragraph{خط دوم (شدید و مزمن)}
\begin{itemize}
\item آسیکلوویر/والاسیکلوویر + سیدوفوویر یا ایمی‌کیمود.
\end{itemize}
\subsection*{ریکارنت هرپس ژنیتالیس}
\begin{itemize}
\item شیوع در زنان بیشتر است.
\end{itemize}
\subsubsection*{تظاهرات بالینی}
\begin{itemize}
\item در \en{HSV2} بین \vih{\en{2--6}} عود در سال قابل انتظار است.
\item کلاسترهایی از وزیکول‌های کوچک.
\item \en{non-indurated ulcers} روی گلنس یا شفت.
\item شدیداً دردناک (در برخی موارد).
\end{itemize}
\subsubsection*{عوارض}
\begin{itemize}
\item فارنژیت: در ۱\% هرپس‌های ژنیتال.
\item احتمال \en{HIV} گرفتن بالا می‌رود.
\end{itemize}
\subsubsection*{سیر بالینی}
\begin{itemize}
\item معمولاً از حملهٔ اولیه کوتاه‌تر است.
\end{itemize}
\subsubsection*{منجمنت}
\begin{itemize}
\item اقدامات تسکینی: کنترل درد؛ پترولیوم ژلی؛ \en{gentle bathing}.
\end{itemize}

\vipage{۱۳}{۱۳}{pic13.jpg — 20261016\_043504613.jpg}
\begin{itemize}
\item فرد بیمار در فاز حاد نباید رابطه داشته باشد.
\item در بین فازها حتماً از کاندوم استفاده کند.
\end{itemize}
\paragraph{خط اول}
\begin{itemize}
\item آسیکلوویر خوراکی \en{ASAP}.
\item \en{Supply of tabs in hands}.
\item رژیم‌های حاشیهٔ تصویر: \en{mg $\times 5$}، \en{mg TDS}، \en{mg BD}؛ همگی \en{5 days}. \viunclear{اعداد ابتدای دوزها خارج از کادر عکس‌اند}.
\item در نقص ایمنی: آسیکلوویر خوراکی \en{BD 5--10 days}؛ \viunclear{عدد دوز در حاشیهٔ چپ بریده است}.
\end{itemize}
\paragraph{خط دوم (عود متعدد / اریتم مولتی‌فرم)}
\begin{itemize}
\item آسیکلوویر خوراکی \en{400 mg BD} روزانه.
\item والاسیکلوویر \en{250 mg BD} یا \en{1 g daily}.
\item فامسیکلوویر \en{125 mg TDS} یا \en{250 mg BD}.
\item مدت رژیم‌های بالا: \en{6--12} ماه.
\item نقص ایمنی (ساپرس): پنسیکلوویر \en{500 mg BD}.
\end{itemize}
\subsection*{هرپس نوزادان}
\begin{itemize}
\item هرپس در ۲۸ روز اول عموماً انتقال از مادر حین زایمان است.
\item هرپس فعال حین زایمان واژینال ریسک بالای انتقال به نوزاد دارد (\viunclear{عدد حاشیه‌ای بریده؛ بخش قابل رؤیت «۵۰\%»}).
\item \en{PROM} می‌تواند هرپس را به جنین منتقل کند.
\item می‌تواند کشنده باشد اگر درگیری مادر اولیه باشد و مادر ایمنی نداشته باشد که به جنین منتقل کند.
\item سزارین اندیکاسیون دارد.
\item داروی آنتی‌وایرال باید داده شود.
\end{itemize}

\vipage{۱۴}{۱۴}{pic14.jpg — 20261016\_043522872.jpg}
\begin{itemize}
\item عفونت اولیهٔ مادر در ابتدای ۳ ماههٔ سوم ریسک عقب‌ماندگی، مرگ جنین و پره‌مچوریتی دارد.
\item فرم راجعه معمولاً بی‌خطر است.
\item هرپس دهانی اولیه خطری برای جنین ندارد.
\item تورچ و راش یا تب در اوایل بارداری: بررسی از نظر \en{HSV}.
\end{itemize}
\subsubsection*{تظاهرات بالینی}
\begin{itemize}
\item ظرف \en{1--3} هفتهٔ اول تولد.
\item می‌تواند پوست، \en{CNS}، چشم، دهان را درگیر کند.
\item به صورت منتشر به بدن: تب، لتارژی، تشنج، دیسترس تنفسی؛ هپاتواسپلنومگالی + هپاتیت و ترومبوسیتوپنی.
\item اگر درمان نشود ۶۰\% احتمال فوت دارد؛ حتی با درمان هم ۲۰\%.
\item بیماری می‌تواند به هموفاگوسیتیک لنفوهیستیوسیتوز پیشرفت کند.
\item درمان با آسیکلوویر است.
\item عود پوستی \viunclear{عبارت کوتاه پس از «پوستی»} شایع است.
\end{itemize}
\subsubsection*{بررسی‌های مورد نیاز}
\begin{itemize}
\item \en{PCR} سواب + کشت ویروس از پوست، مخاط، چشم، دهان و رکتوم.
\item \en{PCR} مایع \en{CSF} در شک به انسفالیت.
\item \en{PCR} خون در شک به بیماری منتشر.
\end{itemize}

\vipage{۱۵}{۱۵}{pic15.jpg — 20261016\_043527398.jpg}
\subsubsection*{درمان}
\begin{itemize}
\item اگر هرپس اولیهٔ فعال در مادر نزدیک به زایمان دیده شود:
\begin{itemize}
\item سزارین.
\item آسیکلوویر پروفیلاکسی برای نوزاد.
\end{itemize}
\item هرپس اولیهٔ ژنیتال حین بارداری: حتماً درمان؛ آسیکلوویر خوراکی.
\item اگر در سه ماههٔ سوم بارداری اپیزود عود داشتیم: از هفتهٔ ۳۶ تا زایمان پروفیلاکسی آسیکلوویر.
\item مادر سابقه داشته ولی الان هرپس فعال ندارد: مانیتور و بررسی نوزاد از نظر هرپس.
\item اگر هرپس نوزادی تشخیص داده شد یا قوی شک بردیم:
\begin{itemize}
\item درمان اورژنت حتی اگر درگیری محدود به پوست باشد.
\item \en{20 mg/kg Aciclovir TDS for 2--3 weeks (IV)}.
\item \en{followed by oral Aciclovir for 6 months}.
\end{itemize}
\end{itemize}
\subsection*{تلقیح هرپس سیمپلکس \en{(Inoculation)} ـ \en{whitlow}}
\begin{itemize}
\item (هرپس پارونیشیال)؛ (هرپتیک ویتلو).
\item عفونت هرپس اولیه ایجاد می‌کند.
\item افرادی که انگشتان خود را گاز می‌گیرند (ویتلو).
\item دورهٔ کمون \en{5--7} روز است.
\item لنفادنوپاتی ناحیه‌ای.
\item علائم کانستیتوشنال خفیف هستند.
\end{itemize}

\vipage{۱۶}{۱۶}{pic16.jpg — 20261016\_043534248.jpg}
\subsection*{هرپس گلادیاتوروم}
\begin{itemize}
\item وزیکول و پوستول روی پلاک به زمینهٔ قرمز در صورت، اسکالپ و قسمت فوقانی تنه.
\item شبیه \en{Impetigo}.
\item \en{10--14} روز طول دورهٔ بیماری است.
\item عود می‌تواند رخ دهد؛ می‌توان با آسیکلوویر پروفیلاکسی خوراکی قبل از مسابقه از عود جلوگیری کرد.
\end{itemize}
\subsection*{عفونت ویروس \en{VZV}}
\begin{itemize}
\item عامل ۲ بیماری واریسلا (آبله‌مرغان، \en{chickenpox}) و زوستر (زونا، \en{shingles}).
\item عفونت اولیه شامل ویرمی و سپس تظاهرات پوستی منتشر است.
\item در ادامه ویروس وارد گانگلیون‌های عصبی (عموماً حسی) می‌شود.
\item ویروس از طریق ترشحات‌های نازوفارنکس منتقل می‌شود.
\item بیماری از \vih{۲ روز قبل تا ۵ روز} ظهور راش‌ها قابل انتقال است.
\item وزیکول‌ها حاوی تعداد بسیار زیادی ویروس هستند و می‌توانند دراپلت تشکیل دهند.
\item ضایعات کراسته ویروس ندارند.
\end{itemize}

\vipage{۱۷}{۱۷}{pic17.jpg — 20261016\_043537144.jpg}
\subsection*{واریسلا (آبله‌مرغان)}
\begin{itemize}
\item اپیدمیولوژی: ۵۰\% موارد اولیه زیر ۵ سال و ۸۵\% قبل از بلوغ.
\item مالتی‌نوکلئیتد جاینت سل با میانگین ۱۵ هسته: کاراکتریستیک.
\item در موارد شدید کبد و سایر ارگان‌ها نیز درگیر می‌شوند.
\item در عفونت \en{VZV}، انکلوزیون اینترانوکلئال، کانالیزیشن فوکال وجود دارد. \viunclear{خوانش «کانالیزیشن فوکال»}.
\item \vih{\en{IgG, IgM, IgA} ظرف \en{3--5}} \vih{روز پس از شروع راش}؛ حداکثر میزان \en{2--3} هفته بعد؛ سپس افت.
\item \en{IgG} در تیتر پایین باقی می‌ماند.
\item آنتی‌بادی‌ها اثر محافظتی ناکامل دارند.
\item ایمنی سلولی مهم‌تر است.
\item در زمستان و اوایل بهار شایع‌تر است.
\item \vih{دورهٔ کمون \en{14--17} روز (\en{9--23})}.
\item ممکن است \vih{\en{1--2} روز تب} و ضعف قبل از ضایعات وجود داشته باشد.
\item ابتدا پاپول ایجاد می‌شود سپس وزیکول‌ها به صورت گروه‌های \en{3--5} تایی، \vih{ظرف \en{2--4} روز}.
\item بیشتر در تنه سپس درگیری صورت و اسکالپ؛ موارد اندام‌ها کمتر.
\item الگوی درگیری \vih{سانتری‌پتال} است.
\item درگیری پروگزیمال اندام‌ها بیشتر است.
\item ضایعات در یک محل در \vih{\en{stage}های مختلف} حضور دارند.
\item وزیکول در دهان و کام شایع است و می‌تواند ملتحمه و ژنیتال را هم درگیر کند.
\end{itemize}

\vipage{۱۸}{۱۸}{pic18.jpg — 20261016\_043544039.jpg}
\begin{itemize}
\item درگیری مخاط مقعد می‌تواند زخم‌های دردناک ایجاد کند.
\item \vih{تب} می‌تواند بسیار خفیف باشد یا می‌تواند تا ۴۱ درجه باشد و \vih{\en{4--5} روز} هم طول بکشد.
\item \vih{خارش} می‌تواند بسیار شدید باشد.
\item فرم هموراژیک در افراد نقص ایمنی شایع است.
\item درگیری مادر باردار با آبله‌مرغان:
\begin{itemize}
\item ریسک ۲۵\% \vih{پنومونیت}؛ بیشترین ریسک در \vih{سه ماههٔ سوم}؛ ریسک نارسایی تنفسی.
\item در \vih{سه ماههٔ اول} و دوم خطر ۲\% آسیب جنینی: سندرم واریسلا کانژنیتال.
\item \vih{حین زایمان} اگر به جنین غیرایمن منتقل شود: ریسک ۳۰\% وفات.
\item \vih{اگر مادر ایمن باشد، در نوزاد بی‌خطر} \vih{است عموماً}.
\end{itemize}
\item انسفالیت: با آتاکسی همراه است معمولاً؛ ۱ در هر ۴۰۰۰ بیمار؛ ریکاوری کامل در ۸۰\%.
\item پنومونی واریسلا در بزرگسالان شایع‌تر از کودکان؛ در سیگاری‌ها و زنان باردار بیشتر می‌شود؛ ضایعات کلسفیه ممکن است رخ دهد.
\item عفونت ثانویهٔ باکتریال پوستی شایع است؛ می‌تواند منجر به اسکار و سپتی‌سمی شود.
\end{itemize}

\vipage{۱۹}{۱۹}{pic19.jpg — 20261016\_043550139.jpg}
\begin{itemize}
\item \en{ITP}: ظرف روز ۵ام تا ۱۰ام رخ می‌دهد؛ معمولاً ظرف \en{3--4} ماه رفع می‌شود؛ می‌تواند در زمینهٔ یک واریسلا خفیف هم باشد.
\item \en{SJS}: در واریسلا و حتی واکسیناسیون هم دیده می‌شود؛ اگر علاوه بر وزیکول، بولا هم دیده شد، ممکن است علاوه بر آسیکلوویر نیاز به کورتون سیستمیک هم باشد.
\item اریتم مولتی‌فرم در واریسلا و زوستر دیده می‌شود.
\end{itemize}
\subsubsection*{سیر بالینی واریسلا و پروگنوز}
\begin{itemize}
\item ظرف \vih{۲۴ روز} ضایعات خشک و کراسته شده و در نهایت اگر عفونت ثانویه رخ ندهد \vih{بدون اسکار} رفع می‌شوند. \viunclear{عدد مدت در تصویر احتمالاً «۲۴» یا «۲--۴» است}.
\item بعد از \vih{۴ روز} ضایعهٔ جدید دیگر رخ نمی‌دهد.
\item بهبودی ضایعات \vih{\en{1--2} هفته} طول می‌کشد.
\item هایپو یا هایپرپیگمنتیشن ممکن است چندین هفته طول بکشد.
\item \vih{اسکارهای دپرشن و گرد در ۱۸\% موارد}.
\item ایمنی مادام‌العمر می‌دهد.
\end{itemize}
\subsubsection*{بررسی‌ها}
\begin{itemize}
\item تشخیص اغلب \vih{بالینی} است.
\item توزیع سنتری‌پتال، ضایعات در مراحل مختلف، خارش.
\item \en{PCR} مایع وزیکول و \en{PCR} سواب دهانی.
\item شک به انسفالومیلیت: \en{PCR} مایع \en{CSF}.
\end{itemize}

\vipage{۲۰}{۲۰}{pic20.jpg — 20261016\_043555366.jpg}
\subsubsection*{منجمنت واریسلا}
\begin{itemize}
\item واکسن پره‌اکسپوژر: ۲ دوز با فاصلهٔ ۳ ماه؛ تا ۱۰ سال ایمنی با آنتی‌بادی.
\item ایمونوگلوبولین پست‌اکسپوژر \en{(VZIG)}.
\item پروفیلاکسی ضد ویروسی.
\item \en{VZIG}: ظرف \vih{۱۰ روز} پس از مواجهه می‌توان داد:
\begin{itemize}
\item نوزادان که مادرشان از \vih{۷ روز قبل تا ۷ روز بعد تولد} دچار واریسلا شده است؛ مادر قبل از زایمان \vih{آسیکلوویر وریدی} بگیرد.
\item \vih{نوزادان سالمی} که با واریسلا فعال مواجهه داشته‌اند یا زوستر.
\item برای افراد با ایمنی مختل.
\item کسی که سابقهٔ ابتلا به واریسلا ندارد و ظرف ۳ ماه گذشته حداقل ۱۴ روز کورتیکواستروئید خوراکی مصرف کرده است.
\item زنان باردار غیرایمن.
\end{itemize}
\item در کودکان مواجهه‌یافته با \en{VZV} اگر \vih{واکسن ظرف \en{3--5} روز} از مواجهه داده شود شدت بیماری را تخفیف می‌کند.
\item در فرد با ایمنی طبیعی: \vih{آسیکلوویر خوراکی ۹ روز بعد از} \vih{مواجهه به مدت ۱ هفته}.
\item پروفیلاکسی ضد ویروسی: پردنیزولون ۴۰ میلی‌گرم روزانه برای بیش از ۱ هفته؛ آسیکلوویر پروفیلاکسی نیاز دارد.
\end{itemize}

\vipage{۲۱}{۲۱}{pic21.jpg — 20261016\_043558617.jpg}
\paragraph{خط اول (افراد سالم)}
\begin{itemize}
\item درمان علامتی: \en{Soothing Antiseptic Applications}.
\item اگر عفونت ثانویه: \en{ABT}.
\item در زنان باردار: آسیکلوویر وریدی.
\item عفونت از ۳ هفته قبل زایمان: نوزاد هم باید درمان شود.
\end{itemize}
\paragraph{خط دوم (کودکان)}
\begin{itemize}
\item شروع آسیکلوویر ظرف \en{24 h} اول راش‌ها.
\item اندیکاسیون‌های آسیکلوویر \en{ASAP}: حداکثر \vih{۲ روز از شروع راش}؛ بزرگسالان، بیماری شدید، نوزاد، نقص ایمنی.
\item \en{10 mg/kg TDS or 4 gr daily 5, 7, 10 day}.
\item اگر وریدی می‌دهیم \en{48 h} بعد خوراکی.
\end{itemize}
\subsection*{زوستر}
\begin{itemize}
\item در کودکان و جوانان شایع نیست.
\item از ۵۰ سالگی زیاد می‌شود؛ میانگین سن \vih{۶۵ سال} است.
\item در زنان کمی شایع‌تر از مردان است.
\item تروما، عصب، سرطان، رادیوتراپی، جراحی.
\item ضایعات پوستی شبیه واریسلا است.
\item پارتیکل درون‌سلولی ممکن است برای بیش از ۱ سال باقی بماند، حتی در یک نورالژیا.
\item سطح \en{IgG} حتی از واریسلا نیز بالاتر می‌رود.
\end{itemize}

\vipage{۲۲}{۲۲}{pic22.jpg — 20261016\_043603802.jpg}
\subsubsection*{تظاهرات بالینی زوستر}
\begin{itemize}
\item تظاهر اولیه معمولاً \vih{درد} است.
\item تب، سردرد، خستگی، تندرنس لوکالیزه.
\item معمولاً یک‌طرفه و لوکالیزه است.
\item ممکن است خارش نیز وجود داشته باشد.
\item فاصلهٔ بین درد و ضایعات پوستی:
\begin{itemize}
\item زونای تری‌ژمینال: \vih{۱٫۴ روز}.
\item زونای توراسیک: \vih{۳٫۲ روز}.
\end{itemize}
\item برای \vih{چندین روز} وزیکول‌های جدید ایجاد می‌شوند.
\item لنفادنوپاتی دردناک شایع است.
\item محل‌های شایع درگیری: توراسیک (۵۳\%)، سرویکال (۲۰\%)، تری‌ژمینال (افتالمیک) (۱۵\%)، لومبوساکرال (۱۱\%).
\item می‌تواند تغییرات موقت \en{EKG} با یا بدون \en{CP} ایجاد کند.
\end{itemize}
\subsubsection*{ورژن‌های بالینی}
\begin{itemize}
\item در بارداری به‌ندرت درگیری جنینی هست.
\item زوستر در شیرخوار ممکن است در صورتی که داخل جنینی از واریسلا مادر گرفته باشد.
\end{itemize}

\vipage{۲۳}{۲۳}{pic23.jpg — 20261016\_043608184.jpg}
\subsubsection*{زوستر تری‌ژمینال}
\begin{itemize}
\item در زوستر عصب افتالمیک، در \en{$\tfrac{2}{3}$} بیماران درگیری چشم رخ می‌دهد.
\item وزیکول در اطراف بینی، بین چشم و پل بینی (درگیری عصب نازوسیلیاری؛ \en{Hutchinson sign}).
\item یووئیت، کراتیت، کنژنکتیویت، فلج عضلات اکولار، پروپتوز، اسکلریت (حتی تا ۳ ماه بعد).
\item انسداد شریان رتینال و نکروز آن.
\item درگیری گانگلیون سیلیاری: \en{Argyll--Robertson pupil}.
\end{itemize}
\subsubsection*{هرپس زوستر اوتیکوس (درگیری عصب فاسیال)}
\begin{itemize}
\item درد گوش.
\item \en{facial palsy}: اگر درمان نشود در ۸۰\% موارد دائمی خواهد شد.
\end{itemize}
\subsubsection*{درگیری موتور در زوستر}
\begin{itemize}
\item ضعف عضلانی پروگزیمال دیده می‌شود.
\item ۵۵\% ریکاوری کامل، ۳۰\% بهبودی نسبی.
\item در شکم: به صورت سودوانسداد کولون (\en{Ogilvie's Syndrome}).
\end{itemize}

\vipage{۲۴}{۲۴}{pic24.jpg — 20261016\_043619090.jpg}
\subsubsection*{نورالژی پست‌هرپتیک}
\begin{itemize}
\item شایع‌ترین عارضهٔ زوستر.
\item درد ۱ ماه یا \vih{۳ ماه} پس از شروع زوستر.
\item در \vih{\en{$\tfrac{1}{3}$} بیماران بالای ۶۰ سال}.
\item در درگیری عصب تری‌ژمینال شایع‌تر است.
\item \vih{الگوی درد}: درد سوزشی ادامه‌دار، همراه هایپراستزی؛ اسپاسمودیک نوروژنیک \viunclear{خوانش عبارت این سطر}.
\item \vih{\en{Allodynia}} (درد به دنبال لمس خفیف) شایع است؛ در ۹۰\% موارد.
\item تا ۳ ماه بعد، ریسک واریسلا واسکولوپاتی‌ها بالا است.
\item \vih{اریتم مولتی‌فرم} به دنبال زوستر دیده شده است.
\item \vih{سندرم گیلن‌باره} به دنبال زوستر مشاهده شده است.
\end{itemize}
\subsubsection*{سیر بالینی زوستر و پروگنوز}
\begin{itemize}
\item درد \en{prodromal} پس از ظهور ضایعات رفع می‌شود.
\item بهبودی ظرف \vih{\en{2--3} هفته} در کودکان و \vih{\en{3--4} هفته} در بزرگسالان.
\item زوستر راجعه ممکن است رخ دهد و در ۴۵\% در همان محل است.
\end{itemize}
\subsubsection*{بررسی‌ها}
\begin{itemize}
\item تشخیص \vih{بالینی} است.
\item تنها \en{DDX}، زوستری‌فرم هرپس سیمپلکس است.
\item \vih{\en{PCR}} نشان‌دهندهٔ \en{VZV DNA}؛ تا ماه‌ها مثبت.
\end{itemize}

\vipage{۲۵}{۲۵}{pic25.jpg — 20261016\_043625632.jpg}
\subsubsection*{منجمنت زوستر}
\begin{itemize}
\item ضد درد و درمان‌های موضعی علامتی.
\item در صورت عفونت ثانویه: \en{ABT}.
\item واکسن زوستر: همان واکسن واریسلا ولی دوز بالا؛ \vih{۱ نوبت} داده می‌شود.
\item بیماران ایمونوساپرس: حداقل ۲ هفته قبل از سرکوب ایمنی.
\item واکسن جدید، \vih{بالای ۵۰ سال} توصیه می‌شود.
\item آنتی‌وایرال \en{ASAP} (بهتر \vih{\en{1--2} روز اول}): درگیری دردناک خصوصاً در پیرسنی؛ نقص ایمنی.
\item \en{Aciclovir 800 mg $\times 5$, 7--10 days}.
\item \en{Valaciclovir 1 gr TDS 7 days; Famciclovir 500 TDS 7 days}.
\end{itemize}
\subsubsection*{درمان پست‌هرپتیک نورالژیا}
\begin{itemize}
\item بهتر است \en{Opioid} ندهیم.
\item \vih{\en{TCA}} (آمی‌تریپتیلین، نورتریپتیلین): \en{25 mg daily for 3--6 months}.
\item یادداشت حاشیه‌ای: هایپرتروفی پروستات \viunclear{آغاز واژهٔ نخست در لبهٔ تصویر بریده است}.
\item معایب: \viunclear{واژهٔ نخست پس از «معایب»}، \vih{چاقی}.
\item گاباپنتین: که خوبه.
\item درمان موضعی با کپسایسین.
\end{itemize}

\vipage{۲۶}{۲۶}{pic26.jpg — 20261016\_043630874.jpg}
\subsection*{ویروس \en{EBV}}
\begin{itemize}
\item تماس با بزاق، بوسه، مقاربت.
\item عفونت اولیه \vih{در کودکان بدون علامت} است.
\item در نوجوانان در \en{$\tfrac{2}{3}$} موارد تبدیل به \vih{\en{IMN}} می‌شود.
\end{itemize}
\subsection*{\en{IMN}}
\begin{itemize}
\item یک \vih{بیماری حاد تب‌دار} است + \vih{گلودرد} و \vih{لنفادنوپاتی سرویکال}.
\item دورهٔ کمون حدود \vih{۶ هفته}.
\item \vih{فارنژوتونسیلیت اگزوداتیو} وجود دارد.
\item \vih{هپاتواسپلنومگالی} در ۵۰\% افراد.
\item پتشی در محل برخورد کام سخت و نرم: \en{2--3} روز پس از تب.
\item اگزانتم در ۱۰\% موارد؛ معمولاً \vih{روز \en{4--6}} بیماری.
\item ماکولوپاپولر هستند و پروگزیمال اندام فوقانی، گاهی هم تنه.
\item ادم دور چشم خصوصاً در \en{F}.
\item ظرف \vih{حداکثر ۱ هفته رفع} می‌شود.
\item اگر آمپی‌سیلین یا آموکسی‌سیلین بدهیم، در ۹۰\% موارد پس از \en{7--10} روز بعد از شروع بیماری، ضایعات ماکولوپاپولر ایجاد می‌شود.
\end{itemize}

\vipage{۲۷}{۲۷}{pic27.jpg — 20261016\_043633806.jpg}
\begin{itemize}
\item میالژی‌ها نادر است.
\item \en{EBV} در همراهی با اختلالات خودایمنی (\en{RA, SLE}، شوگرن \viunclear{انتهای حاشیهٔ چپ}).
\item \en{ITP}: تعداد زیر ۱۰۰ می‌آید.
\item نوریت، میوکاردیت، آتاکسی، پارالیزی فاسیال، نوتروپنی.
\item اگر انسفالیت رخ دهد کشنده است.
\end{itemize}
\subsubsection*{سیر بیماری و پروگنوز}
\begin{itemize}
\item بیماری \vih{ظرف ۱ ماه بهبودی} پیدا می‌کند.
\item لنفادنوپاتی گردنی ممکن است تا \vih{۳ ماه} باقی بماند.
\item ضعف هم ممکن است باقی بماند.
\end{itemize}
\subsubsection*{بررسی‌ها}
\begin{itemize}
\item در فاز حاد \vih{لنفوسیتوز با ۱۰\% سلول آتیپیک} وجود دارد.
\item \vih{\en{LFT} به طور شایع مختل} است؛ در ۴\% منجر به \vih{یرقان} می‌شود.
\item آنتی‌بادی‌ها:
\begin{itemize}
\item در ۸۰\% بیماران حاد، آزمایش‌ها مثبت.
\item \en{Early}: \vih{\en{EA, VCA}}؛ \en{IgM} برای هفته‌های \viunclear{عدد حاشیه‌ای «۲ ...»}؛ \en{IgG} برای ماه‌ها مثبت.
\item \en{late}: \vih{\en{EBNA}}.
\end{itemize}
\item \vih{هتروفیل آنتی‌بادی} در ۹۰\% بیماران ظرف \en{1--2} هفته مثبت.
\end{itemize}
\subsubsection*{منجمنت}
\begin{itemize}
\item درمان علامتی؛ درمان حمایتی.
\item درمان ضد ویروس خوراکی بر مونونوکلئوز حاد بی‌تأثیر است.
\item تظاهرات پوستی: رطوبت‌دهنده؛ استروئید موضعی \vih{\en{moderate potent}}؛ بدون درمان ظرف ۱ هفته.
\end{itemize}

\vipage{۲۸}{۲۸}{pic28.jpg — 20261016\_043638611.jpg}
\subsection*{روزئولا اینفانتوم (بیماری ششم) (اگزانتم سوبیتوم)}
\begin{itemize}
\item یک \vih{بیماری حاد تب‌دار} + \vih{راش ماکولوپاپولر}.
\item شایع‌ترین تب اگزانتماتیک در زیر ۲ سال است؛ حداکثر ۳ سالگی.
\item ناشی از \vih{\en{HHV} تیپ ۶ و ۷}.
\item شایع‌ترین سن \en{6--9} ماهگی است.
\item \vih{دورهٔ کمون \en{10--15} روز} است.
\end{itemize}
\subsubsection*{تظاهر بالینی}
\begin{itemize}
\item \vih{شروع آنی تب}: تا ۴۰ درجه؛ \vih{\en{3--5} روز طول می‌کشد}.
\item \vih{معمولاً هیچ علامت همراهی ندارد}.
\item بی‌قراری، التهاب \en{TM}، تورم کام خلفی و اوولا، هایپرمی حلق و ادم پری‌اوربیت.
\item به دنبال قطع تب، \vih{ماکولوپاپول \en{rose-pink} منتشر} در گردن و تنه؛ در ادامه صورت و اندام‌ها گسترش می‌یابد.
\item به‌ندرت وزیکولر می‌شوند.
\item \vih{لنفادنوپاتی} ساب‌اکسیپیتال و خلفی گردن شایع است.
\item انسفالیت آن نادر است.
\item \en{ITP, SJS} ممکن است رخ دهد.
\end{itemize}
\subsubsection*{سیر بالینی و پروگنوز}
\begin{itemize}
\item \vih{راش‌ها ظرف \en{1--2} روز} بدون اسکار یا پیگمانتاسیون رفع می‌شوند.
\item بعضاً عود رخ می‌دهد.
\end{itemize}

\vipage{۲۹}{۲۹}{pic29.jpg — 20261016\_043641586.jpg}
\subsubsection*{بررسی‌های روزئولا}
\begin{itemize}
\item تشخیص: نبود علائم همراه با تب؛ ماکولوپاپول به دنبال رفع تب.
\item مشاهدهٔ \vih{سروکانورژن یا تیتر \en{Ab}}.
\item \en{IgM} \vih{\en{5--7} روز پس از راش} بالا می‌رود؛ حداکثر آن ۲ هفته؛ تا ۲ ماه هم باقی می‌ماند.
\item تست \en{Ab} فقط برای \en{HHV-6} است.
\item ظرف ۲ روز اول بیماری \vih{لکوسیتوز} وجود دارد ولی با شروع راش‌های پوستی، \vih{لکوپنی} رخ می‌دهد + لنفوسیتوز نسبی.
\end{itemize}
\subsubsection*{منجمنت}
\begin{itemize}
\item تنها درمان مورد نیاز، درمان علامتی است.
\item بیماری شدید: \vih{گانسیکلوویر}، والگانسیکلوویر، سیدوفوویر، فوسکارنت.
\end{itemize}
\subsection*{اگزما هرپتیکوم}
\begin{itemize}
\item عفونت منتشر پوست با یک ویروس که معمولاً درگیری وزیکولار خفیف و موضعی می‌دهد.
\item در بیماری با بیماری قبلی پوستی رخ می‌دهد.
\item شایع‌ترین ویروس \vih{\en{HSV-1}} است.
\item ۳\% بیماران درماتیت آتوپیک به هرپس مبتلا می‌شوند.
\item شایع‌ترین سن \vih{دههٔ دوم و سوم} است.
\end{itemize}

\vipage{۳۰}{۳۰}{pic30.jpg — 20261016\_043647108.jpg}
\begin{itemize}
\item شایع‌ترین بیماری پوستی، \vih{درماتیت آتوپیک} است.
\item می‌تواند در زمینهٔ عفونت اولیه یا عود باشد.
\item \vih{تنها در پوستی با درماتیت قابل} \vih{مشاهده رخ می‌دهد}.
\item بیماران اغلب ایمنی طبیعی دارند.
\item ارتباطی با مصرف سیستمیک یا موضعی \en{CS} ندارد، ولی با تاکرولیموس دارد.
\end{itemize}
\subsubsection*{تظاهرات بالینی}
\begin{itemize}
\item اگر در زمینهٔ عفونت اولیه باشد \vih{دورهٔ کمون \en{5--19} روز} (سه هفته).
\item وزیکول‌های ناف‌دار سریعاً تبدیل به پوستول می‌شوند.
\item گاهی محدود به پوست بیمار، گاهی فراتر از آن.
\item ممکن است به صورت هموراژیک یا تاولی شوند و منطقهٔ بیمار وسیعی ایجاد کنند.
\item پوست قرمز و دردناک است.
\item وزیکول‌ها \en{2--4 mm} قطر دارند و به سرعت پاره می‌شوند.
\item \vih{\en{5--7} روز وزیکول‌های جدید ایجاد می‌شوند}.
\item \vih{تب شدید \en{2--3} روز پس از شروع} \vih{تظاهرات پوستی آغاز می‌شود}.
\item \vih{لنفادنوپاتی موضعی} شایع است.
\item \vih{توکسمی} خصوصاً در عفونت اولیه شایع است.
\item مهم‌ترین \en{DDX} آن، آبله‌مرغان است.
\end{itemize}

\vipage{۳۱}{۳۱}{pic31.jpg — 20261016\_043651337.jpg}
\subsubsection*{سیر بیماری و پروگنوز اگزما هرپتیکوم}
\begin{itemize}
\item \vih{تب ظرف \en{4--5} روز قطع می‌شود}.
\item پوستول‌ها آهسته خوب می‌شوند؛ \vih{اسکار خفیف} \viunclear{کلمهٔ انتهای سطر در حاشیه بریده است}.
\item به‌ندرت انسفالیت رخ می‌دهد.
\item عفونت باکتریال ثانویه می‌تواند سپسیس ایجاد کند.
\item احتمال عود وجود دارد (تا ۲۵\%).
\end{itemize}
\subsubsection*{بررسی‌ها}
\begin{itemize}
\item مایع وزیکول از نظر \en{HSV} باید بررسی شود.
\item \vih{بررسی افزایش \en{Ab}} با انتروال سرم در زمان شروع و ۱۰ روز بعد.
\end{itemize}
\subsubsection*{منجمنت}
\begin{itemize}
\item \vih{همهٔ بیماران باید ضد ویروسی درمان شوند}؛ \en{ASAP}: آسیکلوویر، والاسیکلوویر، فامسیکلوویر.
\item فرم شدید نیازمند داروی وریدی است.
\item قطع داروهای سرکوب ایمنی اگر مقدور شود.
\item در بارداری با آسیکلوویر درمان می‌کنیم.
\item ضایعهٔ اطراف چشم: ویزیت چشم‌پزشک.
\item بررسی از نظر عفونت ملتحمه.
\item عفونت ثانویه: \en{ABT}.
\item وقتی آنتی‌وایرال می‌دهیم می‌توانیم \en{TCS} هم بدهیم.
\item عود مکرر: پروفیلاکسی \vih{\en{long duration}}، آسیکلوویر یا والاسیکلوویر \viunclear{ادامهٔ حاشیه بریده است}.
\item \en{Oral Aciclovir 400 mg $\times 5$ for 5 days}؛ \viunclear{یادداشت حاشیه‌ای آغاز سطر بریده است}.
\end{itemize}

\vipage{۳۲}{۳۲}{pic32.jpg — 20261016\_043656612.jpg}
\subsection*{اورف \en{(Orf)} (اکتیما کنتاژیوزوم)}
\begin{itemize}
\item ویروس \en{ORFV} از خانوادهٔ پاکس‌ویروس‌ها، پاراپاکس‌ویروس‌ها.
\item معمولاً به صورت یک ضایعهٔ منفرد پس از تماس با حیوان آلوده (گوسفند و بز).
\item در گوشت منجمد هم دیده می‌شود؛ در کارکنان.
\end{itemize}
\subsubsection*{تظاهرات بالینی}
\begin{itemize}
\item ظرف \vih{\en{3--10} روز (\en{5--6}) پس از مواجهه}.
\item پاپول منفرد، ماکولر، سفت، قرمز در محل برخورد.
\item تبدیل به یک وزیکول یا بولای هموراژیک، تا پوستول می‌شود.
\item سایز \en{2--3} سانتی‌متر تا ۵ سانتی‌متر.
\item لنفادنیت، لنفانژیت، تب خفیف.
\item می‌تواند مثل آبله‌مرغان منتشر هم بروز کند.
\end{itemize}
\subsubsection*{عوارض}
\begin{itemize}
\item اریتم مولتی‌فرم عارضهٔ شایع است؛ \en{2--3} هفته بعد از اورف شروع می‌شود، در ۱۰\% بیماران.
\end{itemize}
\subsubsection*{سیر بالینی}
\begin{itemize}
\item \vih{بهبودی ضایعه طی \en{3--6} هفته}.
\item ایمنی ایجاد می‌شود و \vih{ضایعهٔ ثانویه نادر است}.
\item سروولوژی تا \vih{۱ هفته پس از ضایعه} قابل بررسی است.
\end{itemize}
\subsubsection*{منجمنت}
\begin{itemize}
\item درمان خاصی نیاز ندارد.
\item اگر عفونت ثانویه: \en{ABT}.
\item ضایعهٔ بزرگ: جراحی.
\end{itemize}

\vipage{۳۳}{۳۳}{pic33.jpg — 20261016\_043700019.jpg}
\begin{itemize}
\item ضایعهٔ مکرر یا ماندگار: کرایوتراپی، ایمی‌کیمود، ایدوکسوریدین، سیدوفوویر.
\end{itemize}
\subsection*{مولوسکوم کنتاژیوزوم \en{(MC)}}
\begin{itemize}
\item شیوع \en{5--11}\%.
\item ویروس \en{MCV}: بیشتر کودکان؛ مخصوص انسان.
\item ۹۰\% به دنبال \vih{\en{MCV-1}}؛ \en{MCV-2} در بزرگسالان شایع‌تر.
\item تظاهر بالینی هر دو تیپ یکسان است.
\item به دنبال تماس با پوست یا اشیای آلوده منتقل می‌شود.
\item زیر ۱ سال نادر است؛ \vih{پیک بین \en{2--5} سال} است.
\item استخر شنا، دستشویی باعث انتقال می‌شوند.
\item در نوجوانان به دنبال ارتباط جنسی، ناحیهٔ ژنیتال شایع است.
\item در پیوند عضو احتمال شدت بیماری افزایش می‌یابد.
\end{itemize}
\subsubsection*{تظاهر بالینی}
\begin{itemize}
\item \vih{دورهٔ کمون ۱۴ روز تا ۶ ماه} است.
\item \vih{پاپول‌های همیسفریک درخشان با ناف مرکزی}.
\item تدریجاً بزرگ می‌شوند؛ ظرف \en{6--12} هفته تا \en{5--10 mm} می‌رسند.
\item به دنبال چسبیدن چند پاپول: پلاک.
\item به طور خودبه‌خودی \vih{پس از چندین ماه} رفع می‌شوند.
\item در نواحی \en{temperate}: گردن، تنه، ران، آگزیلا؛ نواحی \en{Tropic}: اندام‌ها.
\item ضایعات در مراحل مختلف.
\end{itemize}

\vipage{۳۴}{۳۴}{pic34.jpg — 20261016\_043705401.jpg}
\begin{itemize}
\item ضایعه روی لب، زبان، مخاط بوکال: \vih{نقص ایمنی}؛ حتی در ملتحمه بدون ضایعه روی پوست.
\end{itemize}
\subsubsection*{عوارض مولوسکوم}
\begin{itemize}
\item واکنش اگزاماتوز \en{patchy} در ۱۰\% موارد.
\item اریتم مولتی‌فرم.
\item کنژنکتیویت.
\end{itemize}
\subsubsection*{سیر بالینی}
\begin{itemize}
\item اغلب موارد \vih{ظرف \en{6--9} ماه} خودمحدودشونده است.
\item ممکن است تا \vih{\en{3--4} سال} طول بکشد.
\item اگر تنها یک ضایعه باشد معمولاً \viunclear{ادامهٔ سطر خط خورده است}.
\item بیش از ۲ ماه می‌ماند و احتمال تا ۵ سال هست.
\item \vih{\en{BOTE Sign}} (آغاز پایان): نشان می‌دهد ظرف چند ماه آینده رفع خواهد شد؛ التهاب اطراف ضایعات: نشانهٔ عفونت نیست.
\item اسکار دپرشن برجا نمی‌گذارد.
\item \vih{تشخیص بالینی است}.
\end{itemize}
\subsubsection*{منجمنت}
\begin{itemize}
\item در اکثر موارد درمان ضروری نیست.
\item \vih{نباید ضایعات خارانده شوند}.
\item اگزما همراه با مرطوب‌کننده و یک \en{TCS} ضعیف.
\item اگر رفع ضایعات نشد درمان می‌کنیم.
\end{itemize}

\vipage{۳۵}{۳۵}{pic35.jpg — 20261016\_043708970.jpg}
\begin{itemize}
\item داروهای موضعی محرک: رتینوئیک / پودوفیلوتوکسین.
\item \en{Topical Cantharidin 0.7\% w/v q3week max 4 application}.
\item داروهای مدولاتور ایمنی: دیفن‌سیپرون؛ ایمی‌کیمود \viunclear{نام دوم در اصل خط خورده است}.
\item درمان فیزیکی: کرایوتراپی (هر \en{3--4} هفته).
\item ضد ویروسی: \vih{سیدوفوویر موضعی} \en{1--3}\% یا کرم، برای ضایعات شدید.
\end{itemize}
\subsection*{اریتم انفکتیوزوم (پاروویروس \en{B19}) (بیماری پنجم)}
\begin{itemize}
\item در \vih{کودکان} بسیار شایع است.
\item انتقال از طریق \vih{دراپلت}، از طریق جفت.
\end{itemize}
\subsubsection*{تظاهر بالینی}
\begin{itemize}
\item در کودکان \vih{اگزانتم به طور ناگهانی، بدون پرودروم} ظاهر می‌شود؛ بعضاً تب، سردرد، ضعف وجود دارد و تهوع.
\item \vih{پاپول‌های \en{rose-red}} روی گونه‌ها به سرعت به هم متصل شده و یک ظاهر اریتم بالا آمده ایجاد می‌کنند: \vih{\en{slapped cheeks}}.
\item معمولاً \vih{رنگ‌پریدگی پری‌اورال} وجود دارد.
\item ظرف \vih{\en{2--4} روز بعد}، ماکولوپاپول در پروگزیمال اندام‌ها ایجاد و به دیستال حرکت می‌کند؛ به سمت تنه.
\end{itemize}

\vipage{۳۶}{۳۶}{pic36.jpg — 20261016\_043715896.jpg}
\begin{itemize}
\item کف دست و پا ممکن است درگیر شود.
\item ممکن است ماکول‌های قرمز تیره در مخاط بوکال و ژنیتال دیده شوند.
\item \vih{لنفادنوپاتی شایع نیست}.
\item در بزرگسالان: \en{$F<M$} است.
\item پرودروم به صورت \vih{علائم آنفلوانزایی} شایع است.
\item قرمزی صورت معمول نیست.
\end{itemize}
\subsubsection*{عوارض}
\begin{itemize}
\item در مراحل اولیه \vih{لکوسیتوز با لنفوپنی نسبی} قابل انتظار است.
\item در ادامه: \vih{ائوزینوفیلی + لنفوسیتوز}.
\item \vih{آنمی خفیف} شایع است؛ در موارد نادر \en{AA}.
\item میوکاردیت، هپاتیت شدید در موارد اندک.
\item در \vih{۳ ماههٔ اول بارداری}: مرگ جنین، هیدروپس فتالیس.
\end{itemize}
\subsubsection*{سیر بالینی}
\begin{itemize}
\item قرمزی صورت ظرف چند روز رفع می‌شود.
\item ضایعات پوستی ظرف \en{6--10} روز رفع می‌شوند.
\item آرتروپاتی ظرف چند هفته رفع می‌شود؛ در ۱۰\% زنان بیشتر از ۲ ماه.
\item بررسی \en{RA, SLE}.
\end{itemize}

\vipage{۳۷}{۳۷}{pic37.jpg — 20261016\_043719047.jpg}
\subsubsection*{بررسی‌های اریتم انفکتیوزوم}
\begin{itemize}
\item تأیید تشخیص با سرولوژی است.
\item \en{anti VP1 Ab (IgM), anti VP2 Ab (IgM)}: \vih{از هفتهٔ دوم} مثبت می‌شوند؛ \vih{برای ۲ ماه} مثبت می‌مانند.
\item \vih{\en{IgG} تا مدت‌ها می‌ماند}.
\item \en{DNA} ویروس در \vih{\en{PCR} سرم}: سریع‌ترین راه تشخیص.
\end{itemize}
\subsubsection*{منجمنت}
\begin{itemize}
\item \vih{معمولاً درمان نیاز ندارد}.
\item در زنان باردار تشخیص سریع ضروری است.
\item در سرکوب ایمنی، \en{SCD}، بیماری هماتولوژیک: بررسی از نظر نارسایی مغز استخوان.
\end{itemize}
\subsection*{سرخجه (روبلا) (سرخک آلمانی)}
\begin{itemize}
\item عمدهٔ هم آن تأثیر بر جنین است.
\item \vih{یکی از اگزانتم‌های شایع دوران کودکی} \vih{و نوجوانی است}.
\item در بهار شایع‌تر است.
\item انتقال از طریق \vih{دراپلت}.
\item \vih{تا ۴ روز پس از بروز راش} قابل انتقال است.
\end{itemize}
\subsubsection*{تظاهرات بالینی}
\begin{itemize}
\item دورهٔ کمون حدود \vih{۱۴ روز} (\en{12--21}).
\item \vih{کودکان معمولاً پرودروم ندارند}.
\end{itemize}

\vipage{۳۸}{۳۸}{pic38.jpg — 20261016\_043725356.jpg}
\begin{itemize}
\item در بزرگسالان پرودروم \vih{\en{1--5} روز} طول می‌کشد و تب ۳۹ درجه، سردرد، خستگی، گلودرد \vih{بدون کوریزا} دارد.
\item احساس شن در چشم، مثل کنژنکتیویت: \en{gritty sensation}.
\item \vih{به دنبال شروع راش، علائم رفع می‌شوند}.
\item \vih{\en{Forchheimer Sign}}: \vih{انانتم} به صورت ماکول \en{dull red} یا پتشی محدود به \vih{کام نرم} در ۲۰\% بیماران، طی پرودروم.
\item \vih{لنفادنوپاتی} از \en{5--7} روز قبل از راش؛ ژنرالیزه، بیشتر \vih{ساب‌اکسیپیتال}، پست‌اوریکولار و سرویکال.
\item \en{1--2} روز اول تندرنس دارند و برای چند هفته می‌مانند.
\item \vih{راش از صورت شروع شده و به سمت تنه و} \vih{اندام می‌رود}: ماکول‌ها صورتی.
\item \vih{ظرف روز دوم} راش از صورت شروع به پاک شدن می‌کند.
\item \en{WBC} معمولاً \en{NL} ولی می‌تواند \vih{لکوپنی + افزایش پلاسماسل}.
\end{itemize}
\subsubsection*{عوارض}
\begin{itemize}
\item در نوجوانان و بزرگسالان در ۷۰\% زنان و ۸\% مردان \vih{آرتریت} رخ می‌دهد.
\item آرتریت معمولاً ظرف \vih{چند ماه} رفع می‌شود.
\item پورپورا با \vih{\en{ITP}} یا بدون آن.
\item انسفالیت نادر است.
\end{itemize}
\subsubsection*{در بارداری}
\begin{itemize}
\item عفونت اولیهٔ مادر \vih{طی ۱۱ هفتهٔ نخست}: ریسک ۸۵\% آسیب روی جنین؛ آسیب‌های مختلف.
\end{itemize}

\vipage{۳۹}{۳۹}{pic39.jpg — 20261016\_043728010.jpg}
\begin{itemize}
\item طی هفته‌های ۱۲ و ۱۶: ۳۵\%؛ شایع‌ترین آسیب ناشنوایی.
\item آسیب‌های سرخجهٔ کانژنیتال عموماً نورولوژیک هستند.
\item عقب‌ماندگی ذهنی عموماً تا سال اول شناسایی نمی‌شود.
\item میکروسفالی.
\item \en{ITP} نوزادی، زردی نوزادی، آسیب‌های استخوانی.
\end{itemize}
\subsubsection*{سیر بیماری و پروگنوز}
\begin{itemize}
\item ظرف روز سوم راش‌ها تنه رفع می‌شوند.
\item ظرف روز چهارم راش‌ها اندام‌ها رفع می‌شوند.
\item ۴۰\% موارد سرخجه بدون راش هستند.
\end{itemize}
\subsubsection*{بررسی‌ها}
\begin{itemize}
\item \vih{تأیید تشخیص در زنان باردار و} \vih{بیماران نقص ایمنی ضروری است}.
\item \vih{یک نمونه گرفته \en{ASAP} پس از راش،} \vih{دومی \en{7--10} روز بعد}.
\item آنتی‌بادی از زمان ظهور راش قابل شناسایی است (\en{IgM}).
\item \en{IgG}، \en{2--3} هفته بعد قابل شناسایی می‌شود.
\item \en{PCR} سواب بزاق و گلو: سریع‌ترین راه؛ \viunclear{کلمهٔ حاشیهٔ چپ بریده است}.
\end{itemize}
\subsubsection*{منجمنت}
\begin{itemize}
\item پروفیلاکسی: واکسن \en{MMR}؛ دوز اول \en{1--2} سالگی؛ دوز دوم ۷ سالگی؛ زنان، کادر پزشکی و معلمان.
\item \vih{تا ۴ هفته پس از واکسن نباید باردار شد}.
\item آرترالژی عارضهٔ شایع واکسن است.
\end{itemize}

\vipage{۴۰}{۴۰}{pic40.jpg — 20261016\_043733351.jpg}
\begin{itemize}
\item زنان واکسینه‌شده مبتلا نمی‌شوند.
\item برای سایر شرایط دارو یا ایمونوگلوبولین وجود ندارد.
\end{itemize}
\subsection*{بیماری پا و دهان (آفتوویروس)}
\begin{latin}Foot and mouth disease\end{latin}
\begin{itemize}
\item \vih{بیماری حیوانات مزرعه} است؛ به‌ندرت انسان را مبتلا می‌کند.
\item هم در کودکان و هم بزرگسالان که در تماس مستقیم با دام هستند دیده می‌شود.
\item وزیکول‌های کوچک در عمق اپی‌تلیوم مخاط‌ها، حتی حلق.
\end{itemize}
\subsubsection*{تظاهرات بالینی}
\begin{itemize}
\item \vih{دورهٔ کمون \en{2--18} روز}.
\item خستگی، سردرد، تب، \vih{سوزش مخاط دهان}: \en{2--3} روز.
\item سپس وزیکول در مخاط بوکال، زبان و لب‌ها؛ برخی موارد کف دست و پا و پوست بین انگشتان.
\item ممکن است \vih{زخم‌های دردناک و ادماتو} ایجاد کند.
\item بیماری در کودکان و شیرخواران شدیدتر است.
\end{itemize}
\subsection*{عفونت انتروویروس \en{(EV)}}
\begin{itemize}
\item ضایعات ماکولوپاپولر \en{$\pm$} فارنژیت یا درگیری دهان.
\item در تابستان و پاییز شایع‌تر هستند.
\item انتشار سریع فرد به فرد.
\end{itemize}

\vipage{۴۱}{۴۱}{pic41.jpg — 20261016\_043740144.jpg}
\begin{itemize}
\item شایع‌ترین ویروس \en{EV-A} است.
\end{itemize}
\subsubsection*{تظاهر بالینی انتروویروس}
\begin{itemize}
\item \vih{یک بیماری تب‌دار با گاستروانتریت یا} \vih{عفونت تنفسی}.
\item راش‌های کوتاه‌مدت.
\item \vih{راش حتماً پس از تب رخ می‌دهد}؛ معمولاً ماکولوپاپولر یا \viunclear{واژهٔ حاشیه‌ای بریده است}.
\item میوکاردیت، نارسایی تنفسی، مننگوانسفالیت نادر است.
\end{itemize}
\subsubsection*{بررسی‌ها}
\begin{itemize}
\item \en{RT-PCR}: مدفوع، سواب گلو، مایع وزیکول، ترشحات نازال.
\item بسیار مشابه روبلا می‌تواند باشد.
\item \vih{هر زن باردار در اوایل حاملگی که} \vih{دچار راش پوستی می‌شود باید از نظر} \vih{روبلا و \en{EV} حتماً بررسی شود}.
\end{itemize}
\subsection*{بیماری دست، پا، دهان}
\begin{itemize}
\item ناشی از یک \vih{\en{EV}} (معمولاً کوکساکی) می‌باشد.
\item ضعف خفیف، زخم‌های دهانی، وزیکول‌های دست و پا، تب خفیف.
\item \vih{در زمستان شایع است}.
\item بیشتر در \vih{کودکان زیر ۱۰ سال} ولی در بزرگسالان هم شایع است.
\item دورهٔ کمون \vih{\en{5--7} روز}؛ \vih{دورهٔ بیماری یک هفته}.
\item در بزرگسالان استوماتیت دردناک وجود دارد ولی در کودکان نه.
\item ضایعات روی \vih{باسن} هم دیده می‌شوند.
\item در کل بدن معمولاً کمتر از \vih{۵۰ تا}؛ پس از \vih{\en{2--3} روز} از بین می‌روند.
\end{itemize}

\vipage{۴۲}{۴۲}{pic42.jpg — 20261016\_043746349.jpg}
\begin{itemize}
\item \vih{تغییرات ناخن} ممکن است \vih{\en{2--4} هفته} پس از بیماری رخ دهد: اونیکولایزیس، \en{Beau lines}، اونیکومادزیس (افتادن ناخن).
\item شدینگ ویروسی تا ۳ هفته پس از بیماری ادامه دارد.
\item \vih{عود آن نادر است}.
\item درمان خاصی نیست؛ اگر علائم نورولوژیک شدید: \vih{\en{IVIG}}.
\end{itemize}
\subsection*{سرخک \en{(Measles)} (همان روبئولا) (بیماری اول)}
\begin{itemize}
\item ضایعات \vih{ماکولوپاپولر} در فردی با \vih{پرودروم} سردرد، \vih{کوریزا} و کنژنکتیویت؛ \viunclear{عبارت پس از کنژنکتیویت}.
\item زمان شایع \vih{اواخر زمستان و بهار} است.
\item انتقال از طریق تنفس، \en{Semen}، بزاق، مخاط.
\item بیمار از \vih{۴ روز قبل از راش تا ۴ روز بعد از راش} ناقل است.
\item در مناطق اندمیک بیشتر کودکان و در مناطق با پوشش واکسن بیشتر \en{10--19} سال، بزرگسالان جوان.
\item پس از شروع راش‌ها \vih{سرکوب ایمنی \en{T-cell mediated}} رخ می‌دهد و می‌تواند تا \vih{چند ماه} ادامه داشته باشد (\vih{\en{2--6} ماه}).
\item \vih{دورهٔ کمون ۱۰ روز} است.
\end{itemize}
\subsubsection*{تظاهرات بالینی}
\begin{itemize}
\item پرودروم: تب، خستگی، علائم تنفس فوقانی (فارنژیت، رینیت).
\item \vih{فتوفوبی} به دلیل درگیری ملتحمه.
\end{itemize}

\vipage{۴۳}{۴۳}{pic43.jpg — 20261016\_043748551.jpg}
\begin{itemize}
\item از روز دوم \vih{\en{Koplik spots}} روی مخاط بوکال در مقابل دندان‌ها پره‌مولار قابل مشاهده است؛ وجود آن \vih{تشخیصی} است.
\item تب و بی‌اشتهایی برای \vih{\en{3--5} روز} افزایش می‌یابد.
\item اگزانتم ظرف \vih{روز چهارم} روی \vih{پیشانی و پشت گوش‌ها} شروع می‌شود.
\item ظرف \vih{\en{24 h} بعد}، \en{face} را کامل درگیر می‌کند و تنه را درگیر کرده؛ کف دست و پا هم درگیر می‌شود.
\item \vih{راش کمی خارش‌دار است}؛ در ابتدا ماکولار است ولی \en{dull red papules} می‌شود.
\item در موارد شدید راش می‌تواند خونریزی‌دهنده باشد.
\item بیماری که نقص ایمنی سلولی دارد دچار فرم شدید بیماری می‌شود.
\end{itemize}
\subsubsection*{عوارض}
\begin{itemize}
\item برونکوپنومونی، انتریت، اوتیت مدیا.
\item انسفالیت: ۱ در ۱۰۰۰ مورد.
\item \en{SSPE} (پان‌انسفالیت اسکلروزان تحت‌حاد): \en{7--10} سال بعد؛ در بیمار زیر ۱ سال شایع‌تر است.
\item \vih{در بارداری}: ریسک \en{SA} (سقط خودبه‌خودی)، زایمان زودرس \en{(PD)}، بیماری شدیدتر در مادر.
\end{itemize}

\vipage{۴۴}{۴۴}{pic44.jpg — 20261016\_043757184.jpg}
\subsubsection*{سیر و پروگنوز سرخک}
\begin{itemize}
\item \vih{راش از سر و گردن شروع می‌شود و به} \vih{پایین می‌رود}.
\item \vih{از روز ۴ تا ۱۰ راش‌ها رفع می‌شوند}.
\item نقص ایمنی موقت ریسک عفونت باکتریال را بالا می‌برد.
\item ایمنی آن عمرانه است.
\end{itemize}
\subsubsection*{بررسی‌ها}
\begin{itemize}
\item \en{IgM} در سرم و بزاق: \vih{\en{3--4} روز پس از راش}.
\item حداکثر \en{IgG}: \vih{\en{2--4} هفته بعد}.
\item \en{PCR} روتین نیست.
\end{itemize}
\subsubsection*{منجمنت ـ \en{Prevention}}
\begin{itemize}
\item ایمنی غیرفعال، ایمونوگلوبولین انسانی نرمال ظرف \vih{۵ روز} از مواجهه:
\begin{itemize}
\item \vih{مادران باردار غیرایمن} با مواجهه.
\item \vih{کودکان با ریسک خاص}.
\end{itemize}
\item واکسن \en{MMR}: ۳ دوز، (\en{1--2} سالگی)، (۷ سالگی).
\item با وجود واکسن بیماری کم است.
\item ممکن است ظرف ۳ هفته پس از واکسن یک بیماری خفیف رخ دهد.
\item اگر عفونت باکتریال: \en{ABT}.
\end{itemize}

\vipage{۴۵}{۴۵}{pic45.jpg — 20261016\_043800075.jpg}
\begin{itemize}
\item \vih{کودکان} مبتلا به سرخک، \vih{\en{Vit A}} بگیرند: \en{200,000 IU $\times 2$ if $\geq 1$ year}؛ زیر یک سال دوز کمتر.
\item در \vih{سرخک شدید}: ریباویرین.
\end{itemize}
\subsection*{سندرم \en{TORCH}}
\begin{itemize}
\item عفونتی که به صورت کانژنیتال یا پری‌ناتال \viunclear{انتهای جمله در حاشیهٔ چپ بریده است}.
\item توکسوپلاسما، \en{others}، سیفلیس، روبلا، \en{CMV}، \en{HSV}.
\item زردی، پورپورا، تشنج.
\item وزیکول و زخم مخاطی در هرپس.
\item کانسپت \en{TORCH} خیلی بالینی کاربردی نیست.
\end{itemize}
\subsection*{\en{Papular acrodermatitis: Gianotti--Crosti Syndrome (GCS)}}
\begin{itemize}
\item یک \vih{واکنش پوستی} کاراکتریستیک و خودمحدودشونده به عفونت ویروسی.
\item پاپول‌های قرمز روی اندام و صورت.
\item عموماً در \vih{کودکان ۶ ماهه تا ۱۲ ساله}.
\item در بزرگسالان بیشتر در \en{F}.
\item شایع‌ترین ویروس‌ها: \en{HSV}، \vih{\en{EBV}}، \en{CMV}، \en{HBV}.
\item ایکتر: بیماران \vih{\en{hep B}} دارند.
\end{itemize}
\subsubsection*{تظاهر بالینی}
\begin{itemize}
\item دورهٔ کمون \vih{\en{3--4} روز}.
\end{itemize}
\begin{latin}
Pink to red flat-topped papules first on the \vih{Thighs} and \vih{buttocks}, then \vih{extend down} the legs, extensor aspects of the arms, and \vih{finally the face}.
\end{latin}

\vipage{۴۶}{۴۶}{pic46.jpg — 20261016\_043805625.jpg}
\begin{itemize}
\item \vih{تنه، اسکالپ و کف دست و پا} درگیر نمی‌شوند.
\item توزیع ضایعات \vih{غیرقرینه} است.
\item در موارد \en{HBV} خارش وجود ندارد ولی ویروس‌های دیگر ممکن است.
\item \vih{لنفادنوپاتی ژنرالیزه} شایع است؛ بیشتر \vih{آگزیلاری و اینگوئینال}؛ برای چندین ماه می‌ماند.
\item تب خفیف، اسهال، علائم تنفسی ممکن است باشند یا نباشند.
\end{itemize}
\subsubsection*{عوارض}
\begin{itemize}
\item \vih{درگیری کبد} در موارد \en{HBV} شایع است ولی فقط تست‌ها مثبت بوده.
\item هپاتومگالی.
\end{itemize}
\subsubsection*{سیر بالینی}
\begin{itemize}
\item ضایعات ظرف \vih{\en{2--8} هفته} رفع می‌شوند.
\item دسکوامیشن خفیف وجود دارد.
\item هایپو یا هایپرپیگمنتاسیون موقت خصوصاً در پوست‌های تیره.
\end{itemize}
\subsubsection*{بررسی‌ها}
\begin{itemize}
\item \vih{چک \en{LFT}} ضروری است.
\item باید دنبال عفونت عامل بگردیم.
\item \en{CBC}: لکوپنی؛ ندرتاً لکوسیتوز \en{2--15} \viunclear{واحد و پایان عبارت شمارش سلولی}.
\item \en{ESR} بالا نرفته است.
\end{itemize}
\subsubsection*{منجمنت}
\begin{itemize}
\item مرطوب‌کننده.
\item استروئید موضعی.
\item موارد شدید یا راجعه: ریباویرین.
\end{itemize}

\vipage{۴۷}{۴۷}{pic47.jpg — 20261016\_043808488.jpg}
\subsection*{وارت پوستی}
\begin{itemize}
\item در هر سنی دیده می‌شوند ولی در نوزادی و اوایل کودکی نادر هستند.
\item بیشترین شیوع \vih{نوجوانان و جوانان}؛ سپس شیوع طی \en{20--30} سالگی کاهش \viunclear{انتهای حاشیهٔ چپ}.
\item از طریق تماس مستقیم و غیرمستقیم منتقل می‌شوند.
\item وارت پلانتار به طور شایع در استخرها منتقل می‌شود.
\item دورهٔ کمون \vih{چند هفته تا حداکثر ۲ سال} است.
\item وارت \en{Common}: تیپ \en{HPV-2} است.
\item وارت \en{Plane}: تیپ \en{HPV-3} و \en{HPV-10} است.
\item وارت موزائیک: تیپ \en{HPV-2} است.
\item شیوع: \vih{\en{Common $<$ Plantar $<$ Plane $<$}} \vih{فیلی‌فرم}؛ \en{Common} هفتاد درصد؛ \en{Plantar} بیست‌وچهار درصد؛ \en{Plane} سه‌ونیم درصد. % Comparison symbols reproduce those drawn in the source.
\end{itemize}
\subsubsection*{\en{Common warts}}
\begin{itemize}
\item محل شایع پشت دست و انگشتان؛ در کودکان زانو.
\item اگر زیر ناخن باشد تندرنس دارد.
\item اگر ژنیتال را درگیر کند همیشه عفونت جنسی را مطرح نمی‌کند.
\end{itemize}
\subsubsection*{\en{Plane warts}}
\begin{itemize}
\item در صورت، پشت دست، جلوی ساق، خطی گسترش می‌یابد: \vih{\en{Koebner-like}}.
\end{itemize}
\subsubsection*{وارت قصابان}
\begin{itemize}
\item افرادی که تماس شغلی با گوشت دارند، در ۲۳\% آن‌ها دیده می‌شود.
\item از وارت \en{Common} بزرگ‌تر است.
\item عود زیاد.
\item \en{HPV-2, HPV-7}.
\end{itemize}

\vipage{۴۸}{۴۸}{pic48.jpg — 20261016\_043813730.jpg}
\subsubsection*{سیر بالینی وارت}
\begin{itemize}
\item ممکن است هر زمانی به طور خودبه‌خود بهبود یابند.
\item در کودکان \en{clearance} سریع‌تر از بزرگسالان است.
\item بهبودی در \en{plane} معمولاً با \vih{التهاب و خارش درون وارت} همراه است.
\item دپیگمنتیشن ممکن است به دنبال رفع رخ دهد که ظرف چند ماه \en{OK}.
\item تغییرات بدخیم بسیار نادر است و صرفاً در افراد ضعف ایمنی.
\end{itemize}
\subsubsection*{منجمنت}
\begin{itemize}
\item همهٔ وارت‌ها نیاز به درمان ندارند؛ خودبه‌خود رفع می‌شوند.
\item وارت پلانتار باید با چسب پوشیده شود در استخر.
\item \vih{در افراد با نوروپاتی مثل دیابتی‌ها} \vih{یا افرادی که ترمیم زخم آن‌ها مختل} \vih{است بهتر است روش‌های تهاجمی کمتر} \vih{استفاده شود}، ولی \viunclear{عبارت پایانی: «منع نشده»}.
\item \vih{سالیسیلیک اسید} خط اول است: \en{12--26}\% کلودیون یا ژل؛ روزانه تا زمان رفع وارت.
\item \vih{\en{Adhesive plaster}} حاوی \en{SS 40\%} برای پلانتار مفیدند.
\item حتی ممکن است بیشتر از ۳ ماه نیاز باشد.
\item \en{SS + lactic acid}؛ \en{5-FU}؛ \en{Cantharidin}؛ \en{Imiquimod}؛ \en{Cryotherapy}؛ \en{Laser}.
\item \en{Soft paraffin; Zinc Oxide}.
\end{itemize}

\vipage{۴۹}{۴۹}{pic49.jpg — 20261016\_043816582.jpg}
\paragraph{خط دوم وارت پوستی}
\begin{itemize}
\item \vih{کرایوتراپی}: هر ۳ هفته تا ۳ ماه.
\item لیزر می‌تواند باعث پخش وارت صورت شود.
\item دمای ۴۴ درجه به مدت \en{30 min}.
\item جراحی توصیه نمی‌شود: اسکار / عود.
\end{itemize}
\paragraph{خط سوم}
\begin{itemize}
\item ایمی‌کیمود / پودوفیلوتوکسین.
\end{itemize}
\subsection*{وارت آنوژنیتال}
\begin{itemize}
\item سن شایع \vih{\en{20--34} سال} است.
\item در \en{$F=M$} است.
\item دورهٔ کمون \vih{۳ هفته تا ۲۴ ماه}؛ میانگین \en{3--10} ماه.
\item شایع‌ترین \en{HPV-6, HPV-11}.
\item ممکن است خونریزی یا ترشح داشته باشد.
\item حین بارداری تشدید می‌شود.
\item داروهای ضد ویروس بیماری را تشدید می‌کنند: \vih{آنتی‌رترووایرال} \viunclear{ادامهٔ عبارت در حاشیهٔ چپ بریده است}.
\item ریسک سایر \en{STI}ها را بالا می‌برد.
\item ممکن است پس از چند هفته تا چند سال باقی بمانند.
\end{itemize}
\subsubsection*{منجمنت}
\begin{itemize}
\item \vih{واکسن}.
\item \vih{پودوفیلوتوکسین}: کرم \en{0.15\%}؛ محلول \en{0.5\%} موضعی، ۲ بار؛ \viunclear{یادداشت حاشیه‌ای مدت: «۴ هفته ... ۳ ...»؛ بخشی خارج از کادر}.
\item در زنان حامله کنتراندیکه است.
\item \vih{ایمی‌کیمود}: کرم \en{5\%} یا \en{3.75\%}؛ زیر آن‌ها به ترتیب \en{OD} و \en{BD} نوشته شده؛ ۳ بار در هفته، ۱۶ هفته.
\end{itemize}

\ifnum\ViralInfectionsStandalone=1
  \def\ViralInfectionsFinish{\end{document}}
\else
  \let\ViralInfectionsFinish\relax
\fi
\ViralInfectionsFinish
